\textbf{Proof.}
By \(\mathrm{def:bias\text{-}aware\text{-}interval}\), the event \(\tau\in\mathcal I_{h,J_{\mu,n},J_{r,n}}\) is exactly
\[
\left|\widehat\tau^{\mathrm{sp}}_{h,J_{\mu,n},J_{r,n}}-\tau\right|
\le B_X(h,J_{\mu,n},J_{r,n})+R_S(h,J_{\mu,n},J_{r,n})+R_T(m).
\tag{1}
\]
For every \(\mathbb P_{n,m}\in\mathcal M_{n,m}(a,s,t)\), \(\mathrm{thm:ricot\text{-}attainment}\) assigns probability at least \(1-\gamma\) to (1). The bound is uniform in the law, so taking the infimum over the same model class proves the fixed-class coverage display.

The interval is symmetric about \(\widehat\tau^{\mathrm{sp}}_{h,J_{\mu,n},J_{r,n}}\) with half-length \(B_X+R_S+R_T\); hence its realized length is exactly twice that sum. On a design satisfying \(B_X\le b\) and \(\sum_i\Gamma_i^2\le g^2\), the definitions of the radii give
\[
R_S\le\varepsilon^{-1}\sqrt{2\log(4/\gamma)}\,g,
\qquad
R_T(m)=\sqrt{2\log(4/\gamma)/m}.
\]
Substitution proves the final length bound. These are deterministic statements about the realized design and do not entail an expected-length rate or empirical frontier attainment. \(\square\)